\section{Lean formalization architecture}\label{sec:lean}

\subsection{Target repository layout}

The proposed checked development fits naturally under
\path{Derivations/Lean/IPRADerivations/Foundations/TAPIR/}:

\begin{center}
\begin{tabularx}{\textwidth}{lX}
\toprule
Module & Responsibility \\
\midrule
\code{Kernel.lean} & Indexed \code{Plan}, interpreter, structural macros. \\
\code{Schedule.lean} & Sequential schedule and compile-all theorem. \\
\code{Certificates.lean} & Sealing, observer escape, named bounded witnesses. \\
\code{SearchK.lean} & Arity-polymorphic least search and compatibility aliases. \\
\code{Reflection.lean} & Generic-arity bounded aggregators and reflected-formula combinators. \\
\code{Trace.lean} & Generic clocked trace and invariant macro. \\
\code{HeightList.lean} & Conditional bounded enumeration theorem. \\
\code{Examples.lean} & Minimum, finite strategy, and bisection demonstrations. \\
\bottomrule
\end{tabularx}
\end{center}

The archive mirrors this plan in \code{lean/TAPIR/}.  \code{Kernel.lean} and
\code{Schedule.lean} contain no placeholders.  Larger files are explicit source-oriented blueprints
with every remaining milestone marked by \code{sorry}; they are not represented as compiled code.

\subsection{Milestone 0: pin and audit}

Use the pinned commit from the title page.  Before editing, run the existing derivation build and
axiom audit.  Record the versions of Lean, Mathlib, and the repository lake files.  The implementation
must preserve old declarations such as \code{Search.mu} and
\code{BoundedReflection.bounded_choice}; generic-arity replacements should initially be additional
modules or expose compatibility aliases.

\subsection{Milestone 1: elaborate the plan kernel}

The kernel imports only \code{ConciseIPRA.PR}.  Its proof obligations should all be definitional or a
single call to an existing theorem:

\begin{itemize}
\item \code{Plan.weaken} maps to \code{Prf.weaken};
\item \code{Plan.subst} maps to \code{Prf.subst};
\item \code{Plan.cut} maps to \code{Prf.trans};
\item logical constructors map to corresponding \code{Prf} constructors;
\item \code{Plan.qfInd} maps to \code{PRA.qfInd}.
\end{itemize}

The first compile target is:

\begin{lstlisting}[language=Lean]
#check TAPIR.Plan.run
-- expected:
-- Plan Gamma phi -> Gamma |- [PRAAx] phi
\end{lstlisting}

No classical axioms, quotient types, or semantic interpretations should enter this file.

\subsection{Milestone 2: elaborate schedules}

The schedule module depends only on the kernel.  The key function is structurally recursive on the
schedule.  Likely elaboration issues are limited to implicit indices and membership proofs.  A robust
implementation can split helpers:

\begin{enumerate}
\item \code{Plan.hypOfMem : phi in R -> Plan R phi};
\item \code{dischargeAppend} for \(O++R\);
\item \code{Schedule.proveAll};
\item \code{Schedule.head} as a convenience theorem.
\end{enumerate}

Unit examples should construct a two-node schedule where the second node uses the first, then inspect
the type of \code{Plan.run (Schedule.head ...)}.

\subsection{Milestone 3: certificate macros}

Most declarations are wrappers around existing rules and should be placeholder-free:

\begin{lstlisting}[language=Lean]
def Plan.seal
    (hex : Plan Gamma (Fm.ex theta))
    (body : Plan (theta :: Gamma.map Fm.shift) psi.shift) :
    Plan Gamma psi := Plan.exE hex body

theorem namedWitness (f : PTm 2)
    (h : Gamma |- [PRAAx]
      Fm.ex (lt (var 0) (Syn.shift N) /\
        neg (inst2 f (var 0) (Syn.shift x) = zero))) :
    Gamma |- [PRAAx]
      lt (wit f N x) N /\ neg (inst2 f (wit f N x) x = zero) :=
  bounded_choice f h
\end{lstlisting}

The exact notation in the source uses the repository's scoped formula symbols.  Observer macros
should be formulated first for exact equality, then for a user-supplied compatibility formula rather
than hard-coding rational metrics.

\subsection{Milestone 4: generic-arity search}

Implement the body renaming \(\rho\) from \cref{subsec:searchk}.  Recommended helper definitions are:

\begin{itemize}
\item \code{paramsStep : Fin (k+1) -> PTm (k+2)};
\item \code{liftBody (f : PTm (k+1)) : PTm (k+2)};
\item \code{applyK (g : PR (k+1)) (N : PTm n) (xs : Fin k -> PTm n)};
\item substitution lemmas for all three.
\end{itemize}

Define \code{muFK} using \code{recDef}; derive zero and successor equations from
\code{rec_zero} and \code{rec_succ}.  Port the existing invariant proof one branch at a time.  The
preferred regression test is to prove that \(k=1\) specializes to the current \code{Search.mu}
up to a definitional or object equality, and then leave the old API untouched.

\subsection{Milestone 5: reflection packages}

A practical Lean record is:

\begin{lstlisting}[language=Lean]
structure Reflected {n : Nat} (phi : Fm PTm n) where
  bit      : PTm n
  cases    : ([] : List (Fm PTm n)) |- [PRAAx]
               (bit = zero) \/ (bit = one)
  sound    : ([] : List (Fm PTm n)) |- [PRAAx]
               ((bit = one) --> phi)
  complete : ([] : List (Fm PTm n)) |- [PRAAx]
               (phi --> (bit = one))
\end{lstlisting}

For use under hypotheses and arbitrary substitutions, derive methods by applying
\code{IPRA.inst} or \code{Prf.subst} followed by \code{Prf.ofNil}.  Implement atomic equality and
nonzero first, then conjunction and disjunction.  Negation and implication should explicitly use the
bit cases; do not invoke classical tactics.  Bounded quantifiers are added after \code{SearchK} and
generic aggregators compile.

A source-level bounded formula datatype is optional.  The safest first implementation builds
\code{Reflected} packages directly in theorem-specific modules.  A later meta-program can recurse on
a syntax tree and emit the same packages.

\subsection{Milestone 6: traces}

The generic trace module should use \code{recDef} rather than introduce a new recursion encoding.
A useful theorem signature is:

\begin{lstlisting}[language=Lean]
theorem traceInvariant
    (I : Fm PTm (k+1)) (hI : I.QF)
    (h0 : ([] : List (Fm PTm k)) |- [PRAAx]
      I.subst (Syn.inst zero))
    (hS : [I] |- [PRAAx]
      I.subst (Syn.repl (succ (var 0)))) :
    ([] : List (Fm PTm (k+1))) |- [PRAAx] I :=
  IPRA.ind I hI h0 hS
\end{lstlisting}

The theorem is essentially the existing \code{IPRA.ind}; the new module should focus on state
projection and failure-accumulator combinators, not duplicate induction.

\subsection{Milestone 7: conditional complete lists}

Formalize a generic scan over \([0,C)\) whose membership bit is a reflected predicate.  Prove:

\begin{enumerate}
\item scan soundness;
\item membership of every accepted code below \(C\);
\item completeness from a residual formula \(S(z)\to z<C\).
\end{enumerate}

The first application should be a toy polynomial equation with an elementary coefficient/height
bound.  Only after that compiles should theorem packets for difficult Diophantine problems import the
generic module and expose their true height residual.

\subsection{Milestone 8: examples and integration}

Three examples should be required before merging:

\begin{enumerate}
\item finite minimum with a named least minimizer;
\item a local certificate eliminated without any search bound;
\item a three-node schedule ending in a conditional theorem with one residual, followed by a second
      theorem that supplies the residual and closes the result.
\end{enumerate}

Then update \code{IPRADerivations.lean}, the derivation README, and the root unification plan.  Run
all existing tests and the repository's axiom scan.  The final declaration should visibly have the
form
\[
 ([]:\mathsf{List}(\Fm\,\PTm\,0))\vdash_{\PRAAx}\Phi
\]
or the intended finite residual context.

\subsection{Expected formalization risks}

\paragraph{Dependent pattern matching.}
The indexed \code{Plan.run} definition may need explicit implicit arguments or separate equations.
This is a Lean elaboration issue, not a mathematical gap.

\paragraph{Binder simplification.}
Generic reflection will require lemmas normalizing substitutions through shifted body terms.  These
should be isolated as named equalities rather than handled by large brittle \code{simp} calls.

\paragraph{Context order.}
The schedule uses \(O++R\).  Domain nodes may naturally produce \(R++O\) or small subcontexts.  Add
explicit permutation/weakening helpers and choose one canonical order.

\paragraph{Proof-term size.}
Compositional reflected formulas may duplicate bit-case proofs.  Mark common packages as definitions
and let theorem-specific code reuse them.  If necessary, lower large equational blocks to the checked
native clause calculus.

\paragraph{API drift.}
The branch advances quickly.  Keep all upstream declaration names and blob identifiers in the source
audit, and repin before integration.

\subsection{Definition of done}

A module is complete only when:

\begin{enumerate}
\item it elaborates in the pinned repository with no \code{sorry};
\item all public theorem endpoints are ordinary \code{Prf PRAAx};
\item no new axiom appears in the axiom audit;
\item every host algorithm is absent from the trust argument;
\item residual producer formulas are listed in theorem statements and metadata;
\item at least one theorem-specific client imports the module.
\end{enumerate}
